\section*{Hoofdstuk \ref{ch:b2:neurons-synapses} --- Neuronen, actiepotentialen en synapsen}

\begin{solution}{exo:b2:neurons-synapses:1}
De \omterm{def:b2:neurons-synapses:neuron}{dendrieten} ontvangen synapsen; het soma bevat de kern en integreert; de
axonheuvel beslist; het \omterm{def:b2:neurons-synapses:neuron}{axon} geleidt; de uiteinden geven transmitter af.
\omterm{def:b2:neurons-synapses:neuron}{Myeline} is het omgewikkelde membraan van een gliacel, vrijwel zonder
cytoplasma; het isoleert het \omterm{def:b2:neurons-synapses:neuron}{axon}, zodat de impuls tussen de knopen springt
en honderd keer sneller gaat tegen lagere kosten.
\end{solution}

\begin{solution}{exo:b2:neurons-synapses:2}
Depolarisatie tot de drempel; stijgende fase --- spanningsafhankelijke
natriumkanalen gaan open, positieve terugkoppeling; piek dicht bij
$E_{\mathrm{Na}}$ --- de natriumkanalen inactiveren; dalende fase ---
spanningsafhankelijke kaliumkanalen gaan open; nahyperpolarisatie
--- kaliumkanalen nog open, dicht bij
$E_{\mathrm K}$; herstel wanneer ze sluiten. Alles-of-niets omdat de
positieve terugkoppeling ofwel helemaal doorloopt ofwel helemaal niet; in één
richting omdat het zojuist doorlopen membraan refractair is (natriumkanalen
geïnactiveerd).
\end{solution}

\begin{solution}{exo:b2:neurons-synapses:3}
De impuls bereikt het uiteinde; spanningsafhankelijke calciumkanalen gaan
open; calcium zet de versmelting van blaasjes in gang; de transmitter
diffundeert over de spleet; hij bindt postsynaptische receptoren; kanalen
gaan open (of een G-eiwit schakelt); \omterm{def:b2:neurons-synapses:synapse}{postsynaptische potentiaal}; de
transmitter wordt door een enzym of transporter verwijderd; het membraan van
het blaasje wordt teruggewonnen.
\end{solution}

\begin{solution}{exo:b2:neurons-synapses:4}
Elektrisch: geen vertraging, meestal in beide richtingen, geen versterking,
geen remming (er gaat alleen stroom over). Chemisch: een halve milliseconde,
in één richting, versterking (één impuls geeft honderden quanta af), remming
mogelijk (kanalen voor chloride of kalium), en veranderbaar.
\end{solution}

\begin{solution}{exo:b2:neurons-synapses:5}
Uit $E = (61.5/z)\log_{10}(c_{\text{uit}}/c_{\text{in}})$ volgt:
voor kalium \qty{-89}{mV}; voor natrium \qty{+61}{mV}; voor chloride, met $z =
-1$, $-61.5\log_{10} 11 = \qty{-64}{mV}$; calcium, met $z = 2$,
$30.75\times 4.3 = \qty{+132}{mV}$. Bij \qty{-70}{mV}:
kalium gaat naar buiten, natrium komt binnen, chloride gaat licht naar buiten
(het membraan staat onder zijn evenwicht), calcium komt sterk binnen.
\end{solution}

\begin{solution}{exo:b2:neurons-synapses:6}
$20 : 1$: $(20\times(-89) + 61)/21 = \qty{-82}{mV}$, de rusttoestand dicht
bij $E_{\mathrm K}$. $1 : 20$: $(-89 + 20\times 61)/21 =
\qty{+54}{mV}$, de piek, dicht bij $E_{\mathrm{Na}}$. $50 : 1$: $(50\times
(-89) + 61)/51 = \qty{-86}{mV}$, de nahyperpolarisatie, met extra
kaliumkanalen open.
\end{solution}

\begin{solution}{exo:b2:neurons-synapses:7}
$\sqrt{500} = 22$: \qty{22}{m/s}. Om \qty{100}{m/s} te halen, zou een
ongemyeliniseerd \omterm{def:b2:neurons-synapses:neuron}{axon} $d = 10^{4}$ \unit{\micro m} nodig hebben, een
centimeter. Eén meter kost \qty{10}{ms} in de gemyeliniseerde vezel,
\qty{45}{ms} in het \omterm{def:b2:neurons-synapses:neuron}{axon} van de pijlinktvis, een volle seconde in een vezel
van \qty{1}{\micro m}.
\end{solution}

\begin{solution}{exo:b2:neurons-synapses:8}
$m = \ln(200/74) = 1.0$; $P(1) = e^{-1} = 0.37$, $P(2) = 0.18$, $P(3)
= 0.06$. Gemiddelde respons $1.0\times 0.5 = \qty{0.5}{mV}$.
\end{solution}

\begin{solution}{exo:b2:neurons-synapses:9}
$1/\qty{2}{ms} = 500$ impulsen per seconde. De intensiteit van de aanraking
wordt gecodeerd door de frequentie van de impulsen (en door het aantal
\omterm{def:b2:neurons-synapses:neuron}{neuronen} dat wordt gerekruteerd); de impulsen zelf zijn identiek.
\end{solution}

\begin{solution}{exo:b2:neurons-synapses:10}
$\lambda = \sqrt{R_m d/4R_i}$: \qty{0.5}{mm}, \qty{1.6}{mm},
\qty{11}{mm}. Met \omterm{def:b2:neurons-synapses:neuron}{myeline}, \qty{10}{\micro m}: $1.6\times\sqrt{200} =
\qty{22}{mm}$. Een internodium van \qty{1}{mm} verliest slechts $1 - e^{-1/22}
= \qty{4}{\%}$ van het signaal, zodat de volgende knoop gemakkelijk tot de
drempel wordt gebracht. Over een gedemyeliniseerd stuk van \qty{3}{mm} daalt
het signaal tot $e^{-3/1.6} = \qty{15}{\%}$ --- van een impuls van
\qty{100}{mV} ongeveer de drempel van \qty{15}{mV}: de geleiding faalt of wordt
onbetrouwbaar.
\end{solution}

\begin{solution}{exo:b2:neurons-synapses:11}
Tetrodotoxine: geen natriumstroom, geen impuls, geen overdracht.
Tetraethylammonium: geen kaliumstroom, verlengde impuls, meer transmitter
afgegeven per impuls. Geen calcium: normale impulsen, geen afgifte, geen
overdracht. Curare: normale afgifte, receptoren geblokkeerd, geen
eindplaatpotentiaal --- verlamming. Remmer van cholinesterase: acetylcholine
blijft aanwezig, de receptoren blijven open, de spier depolariseert en kan
daarna niet repolariseren --- verlamming door een teveel. Botulinetoxine:
blaasjes kunnen niet versmelten, geen afgifte --- verlamming.
\end{solution}

\begin{solution}{exo:b2:neurons-synapses:12}
De vertraging koopt: overdracht in slechts één richting; versterking (één
enkele impuls geeft honderden quanta af en kan een grotere cel aansturen);
remming, wat een elektrische verbinding niet kan; en plasticiteit, omdat het
aantal quanta en receptoren met het gebruik kan veranderen, en zo leren
circuits. Elektrische synapsen worden gebruikt waar synchronie en snelheid
meer tellen dan rekenwerk: vluchtcircuits, de hartspier, sommige inhiberende
netwerken.
\end{solution}

\begin{solution}{pb:b2:neurons-synapses:1}
\textbf{1.} $E_{\mathrm K} = 61.5\log_{10}(4/140) = \qty{-95}{mV}$;
$E_{\mathrm{Na}} = 61.5\log_{10}(145/12) = \qty{+67}{mV}$.
\textbf{2.} $(25\times(-95) + 67)/26 = \qty{-89}{mV}$.
\textbf{3.} $E_{\mathrm K} = 61.5\log_{10}(8/140) = \qty{-76}{mV}$;
$V = (25\times(-76) + 67)/26 = \qty{-71}{mV}$: dichter bij de drempel,
dus eerst prikkelbaarder --- en, als het aanhoudt, minder, omdat de
natriumkanalen bij het gedepolariseerde niveau inactiveren.
\textbf{4.} De gradiënten lopen in de loop van uren af doordat natrium naar
binnen en kalium naar buiten lekt; de potentiaal drijft naar nul; met de pomp
stilgelegd trekken de niet-doorlaatbare anionen van de cel chloride en water
aan, en de cel zwelt op.
\textbf{5.} Oppervlak per millimeter $\pi\times 15\times 10^{-6}\times
10^{-3} = \qty{4.7e-8}{m^2}$; $C = \qty{4.7e-10}{F}$; $Q = CV =
\qty{4.7e-11}{C}$: $2.9\times 10^{8}$ ionen.
\textbf{6.} Volume per millimeter $\pi(7.5\times 10^{-6})^{2}\times
10^{-3} = \qty{1.8e-13}{m^3} = \qty{1.8e-10}{L}$, met daarin $12\times
10^{-3}\times 1.8\times 10^{-10} = 2.1\times 10^{-12}$ mol, $1.3\times
10^{12}$ ionen: een verandering van $2\times 10^{-4}$.
\textbf{7.} $2.9\times 10^{8}/3 \approx 10^{8}$ ATP per millimeter per
impuls.
\textbf{8.} $0.9/90 = \qty{10}{ms}$ in elke richting.
\textbf{9.} Naakt: $\sqrt{1\times 15\times 10^{-6}/4} = \qty{1.9}{mm}$;
gemyeliniseerd: $\times\sqrt{250} = \qty{31}{mm}$.
\textbf{10.} Naakt: $e^{-1.5/1.9} = 0.45$, \qty{45}{mV}; gemyeliniseerd:
$e^{-1.5/31} = 0.95$, \qty{95}{mV}. Beide liggen boven de drempel van
\qty{15}{mV}; maar het naakte \omterm{def:b2:neurons-synapses:neuron}{axon} moet onderweg zijn hele membraan opladen,
wat het traag maakt, terwijl onder \omterm{def:b2:neurons-synapses:neuron}{myeline} bijna niets verloren gaat of
wordt opgeladen.
\textbf{11.} $\sqrt{15} = \qty{3.9}{m/s}$: \qty{0.23}{s} in elke richting,
bijna een halve seconde voor de heen- en terugweg --- te traag voor een reflex
die een struikeling moet opvangen.
\textbf{12.} $e^{-4/1.9} = 0.12$: \qty{12}{mV} bij de verre knoop, onder
de drempel --- de impuls wordt geblokkeerd en de reflex gaat verloren.
\textbf{13.} Achter de impuls zijn de natriumkanalen een milliseconde of twee
geïnactiveerd, zodat het zojuist doorlopen membraan niet opnieuw kan worden
geprikkeld en de golf alleen vooruit kan.
\textbf{14.} $m = \ln(300/111) = 1.0$.
\textbf{15.} $P(1) = 0.37$, $P(2) = 0.18$, $P(3) = 0.06$: ongeveer 110, 55
en 18 van de 300 pogingen.
\textbf{16.} $1.0\times 0.4 = \qty{0.4}{mV}$.
\textbf{17.} $P(0) = e^{-60} \approx 10^{-26}$: nooit; gemiddelde respons
\qty{24}{mV}, boven de drempel van \qty{15}{mV}: één zo'n \omterm{def:b2:neurons-synapses:synapse}{synaps} doet het
motorische \omterm{def:b2:neurons-synapses:neuron}{neuron} vuren.
\textbf{18.} Twintig synapsen die samen actief zijn, tellen hun stromen op
(ruimtelijke summatie), minder dan lineair naarmate de potentiaal de
omkeerpotentiaal van de kanalen nadert; zelfs bij elk $m = 1$ zouden ze
\qty{8}{mV} geven, en bij normale $m$ ruim meer dan de drempel.
\textbf{19.} Het openen van chloridekanalen bij $E_{\mathrm{Cl}} = V_{\text{rest}}$
voegt geleidbaarheid toe zonder de potentiaal te verplaatsen; de exciterende
stroom loopt nu door een lekkender membraan en geeft een kleinere
depolarisatie --- shuntinhibitie.
\textbf{20.} $10 + 0.6 + 10 + 0.6 + 5 = \qty{26}{ms}$, tegenover
\qty{30}{ms} gemeten.
\textbf{21.} De eigen activering van de receptor in de spierspoel, de
voortplanting van de actiepotentiaal van de \omterm{def:b2:cell-differentiation:myogenesis}{spiervezel} langs de vezel
(meters per seconde over centimeters), en de mechanische latentie voordat er
kracht verschijnt.
\textbf{22.} Elke \omterm{def:b2:neurons-synapses:synapse}{synaps} voegt een halve milliseconde toe en een
schakelneuron voegt een cel toe; maar schakelneuronen maken het mogelijk het
signaal om te keren (de antagonist remmen), met andere te combineren en door
de hersenen te laten moduleren --- een reflex die kan worden bijgestuurd.
\textbf{23.} $90 - 6.2 = \qty{84}{ms}$ geleiding voor \qty{1.8}{m}:
ongeveer \qty{21}{m/s}.
\textbf{24.} De helft van de inwaartse stroom bij elke spanning: de drempel
stijgt, de impuls wordt kleiner en trager, de veiligheidsmarge bij elke
knoop daalt, en de impuls kan uitvallen; de reflex verzwakt of verdwijnt.
\textbf{25.} \omterm{thm:b2:neurons-synapses:chord}{Rustpotentiaal} \qty{-89}{mV}; \qty{10}{ms} in elke richting;
$m = 1.0$; berekende latentie \qty{26}{ms}.
\end{solution}
